\documentclass[10pt]{article}
\usepackage{icmlko}

\icmlmeta{IP 파운데이션 월드모델}{94}{표지를 삼천오백 장 받아 봤다}{처음으로 메타데이터 밖에서 축을 긁었다 --- 그리고 같은 표지가 도메인마다 다른 것을 재고 있었다}

\begin{document}

\twocolumn[
\icmlheader

\begin{abstract}
노트 93이 남은 격차의 대부분(0.256)이 ``축에 없는 정보''라 했고, 기존 데이터로
만들 수 있는 축은 노트 84 · 89 · 90에서 다 시도했다. \textbf{그래서 새로 긁었다.}
지금 축은 전부 메타데이터(수 · 길이 · 가격)인데 \textbf{표지는 보는 것}이고
소비 결정의 첫 관문이다. AniList 에서 표지 대표 색 \textbf{5,015건}을 받고,
다섯 도메인(애니 · 모바일 · 웹툰 · 만화 · 세계애니)에서 실제 표지 \textbf{3,500장}을
내려받아 밝기 · 채도 · 가장자리 밀도 · 색 엔트로피 · 명암 대비를 뽑았다.
\textbf{첫 신호가 좋았다} --- AniList 대표 색의 명도가 만화 $+$0.101, 세계애니
$+$0.186으로 같은 부호였다. \textbf{그런데 실제 이미지의 평균 밝기는 만화
$-$0.092, 세계애니 $-$0.061로 부호가 반대다.} ``표지 색''이 하나의 물리량이
아니었다 --- 대표 색과 평균 밝기는 다른 것을 잰다. 다섯 특징 $\times$ 다섯
도메인 스물다섯 칸에서 \textbf{부호가 일관된 것이 하나도 없다}(밝기: 애니
$+$0.110 · 모바일 $+$0.109 · 만화 $-$0.092). 축으로 넣으면 판정치가
$-$0.003$\sim-$0.010이다. \textbf{다만 팝업은 다섯 설정 모두에서 오른다}
($+$0.003$\sim+$0.007). 둘 다 문턱 아래라 채택하지 않는다. \textbf{표지는 축이
안 된다.}
\end{abstract}
]

\begin{figure*}[t]
\centering
\includegraphics[width=\textwidth]{figs/cov.pdf}
\caption{왼쪽 --- 표지 특징 다섯과 라벨의 상관. 오른쪽 --- 축으로 넣은 결과.}
\end{figure*}

\section{메타데이터 밖으로 나갔다}

지금 축은 전부 세는 것이다 --- 태그 수 · 언어 수 · 화수 · 가격 · 용량.
\textbf{표지는 다르다.} 웹툰 목록에서도 만화 서가에서도 앱스토어에서도 사람이
먼저 보는 것은 그림이고, 팝업의 포토존 · 체험 밀도와도 이어진다.

\begin{table}[t]
\centering\tsize
\caption{긁은 것.}
\begin{tabular}{lr}
\toprule
AniList 표지 대표 색 & 5,015건 \\
내려받은 표지 이미지 & 3,500장 \\
도메인 & 5개 \\
뽑은 특징 & 밝기 · 채도 · 가장자리 · 색 엔트로피 · 명암 \\
\bottomrule
\end{tabular}
\end{table}

계산은 다섯 도메인에 \textbf{똑같이} 걸었다(노트 74의 규약) --- 100픽셀로 줄여
회색조 평균 · HSV 채도 평균 · 기울기 크기 평균 · 5$\times$5$\times$5 색 상자
엔트로피.

\section{첫 신호가 좋았고, 뒤집혔다}

\begin{table}[t]
\centering\tsize
\caption{같은 작품, 두 가지 ``표지 색''.}
\begin{tabular}{lrr}
\toprule
& AniList 대표 색 명도 & 실제 이미지 평균 밝기 \\
\midrule
만화 & \textbf{$+$0.101} & \textbf{$-$0.092} \\
세계애니 & \textbf{$+$0.186} & \textbf{$-$0.061} \\
\bottomrule
\end{tabular}
\end{table}

\claimx{``표지 색''은 하나의 물리량이 아니다. AniList 가 뽑은 \emph{대표} 색의
명도와 이미지 \emph{평균} 밝기가 같은 작품에서 반대 부호를 낸다.}{대표 색이
어떻게 계산되는지 모른다. 채도 높은 강조색을 고르는 방식이면 평균 밝기와
반대로 움직이는 것이 자연스럽다.}

이것 자체가 새 축을 만들 때의 교훈이다 --- \textbf{이름이 같다고 같은 것을
재는 것이 아니다.} 노트 74가 도메인 사이에서 배운 것을 여기서는 한 도메인
안에서 다시 배운다.

\section{부호가 도메인마다 다르다}

\begin{table}[t]
\centering\tsize
\caption{표지 특징과 라벨의 순위 상관. 덮개율 23$\sim$34\%.}
\begin{tabular}{lrrrrr}
\toprule
& 애니 & 모바일 & 웹툰 & 만화 & 세계애니 \\
\midrule
밝기 & $+$0.110 & $+$0.109 & $-$0.015 & $-$0.092 & $-$0.061 \\
채도 & \textbf{$-$0.188} & $+$0.106 & $+$0.079 & $+$0.038 & $+$0.076 \\
가장자리 & $-$0.092 & $+$0.039 & $+$0.096 & $-$0.071 & $-$0.041 \\
색 엔트로피 & \textbf{$-$0.218} & $+$0.066 & $+$0.081 & $-$0.066 & $+$0.067 \\
명암 대비 & $+$0.014 & $-$0.000 & $+$0.017 & $+$0.047 & $+$0.063 \\
\bottomrule
\end{tabular}
\end{table}

\claimx{표지 특징은 도메인을 건너 같은 방향을 안 가리킨다. 스물다섯 칸에서
부호가 일관된 특징이 하나도 없고, 가장 큰 값(애니 색 엔트로피 $-$0.218)이
다른 도메인에서는 $+$0.067이다.}{덮개율이 23$\sim$34\%다 --- 도메인마다 700장씩만
받았고 레코드 ID 순으로 잘랐다. 그 부분집합이 치우쳤을 수 있다.}

\section{축으로 넣으면}

\begin{table}[t]
\centering\tsize
\caption{덮개율이 낮아 중립 대입(노트 85)으로 넣었다.}
\begin{tabular}{lrr}
\toprule
& $\Delta$판정치 & $\Delta$팝업 \\
\midrule
밝기(고유 축) & $-$0.0092 & $+$0.0066 \\
채도(고유 축) & $-$0.0030 & $+$0.0047 \\
색 엔트로피(고유 축) & $-$0.0101 & $+$0.0032 \\
셋 다(고유 축) & $-$0.0070 & $+$0.0045 \\
밝기(공통 축) & $-$0.0104 & $+$0.0066 \\
\bottomrule
\end{tabular}
\end{table}

판정치는 다섯 다 내리고 \textbf{팝업은 다섯 다 오른다}. 둘 다 구간 반폭(0.033)의
3분의 1 아래라 채택하지 않는다.

\begin{quote}\fontsize{9.5}{11.5}\selectfont
팝업이 일관되게 오르는 것이 흥미롭다. 팝업은 표지가 \emph{없는} 도메인이고
축이 하나도 안 붙는데, 다른 도메인에 시각 축이 붙으면 그 도메인들의 인자 공간이
바뀌어 팝업 예측이 좋아진다. \textbf{다섯 번 다 같은 방향이지만 크기가
$+$0.003$\sim+$0.007이라 잡음과 못 가른다.}
\end{quote}

\section{한계}

\begin{itemize}[leftmargin=1.1em,itemsep=0.15em]
\item 도메인마다 700장씩만 받았고 \textbf{레코드 ID 순}으로 잘랐다. 무작위가
      아니라 치우쳤을 수 있다 --- 가장 큰 결함이다.
\item 표지 특징이 다섯뿐이다. 얼굴 유무 · 인물 수 · 글자 비중 같은 것은
      모형이 필요해 안 했다.
\item 팝업 · 아이돌 · 게임 · 도서 · 펀딩에는 표지 축이 없다. 다섯 도메인만
      붙었다.
\item ``대표 색과 평균 밝기가 반대''를 관측했을 뿐 AniList 의 계산법은 모른다.
\item 덮개율 30\%에 중립 대입이라 신호가 3분의 1로 희석된다(노트 85).
\end{itemize}

\section{다음}

\begin{enumerate}[leftmargin=1.3em,itemsep=0.15em]
\item \textbf{무작위로 다시 받는다} --- ID 순 700장이 가장 큰 결함이다. 무작위
      700장으로 상관표를 다시 낸다. 부호가 바뀌면 결론도 바뀐다.
\item 다음 후보는 \textbf{텍스트}다 --- 줄거리 · 소개문. 표지보다 도메인 간
      대응이 분명하고(모든 도메인이 설명문을 갖는다) 뽑을 수 있는 양도 많다.
\item 얼굴 · 인물 수를 뽑으려면 모형이 필요하다. 값어치를 먼저 가늠한다.
\item 팝업이 다섯 번 다 오른 것 --- 다른 도메인의 축이 팝업을 돕는 경로가
      있는지 따로 본다.
\end{enumerate}

\vspace{0.6em}
\hrule height 0.3pt
\vspace{0.5em}
{\footnotesize
\textbf{참고문헌}\\[0.3em]
Grinsztajn, L. et al. (2022). Why do tree-based models still outperform deep
learning on typical tabular data? \emph{NeurIPS}.\\[0.15em]
Little, R. J. A., Rubin, D. B. (2019). \emph{Statistical Analysis with Missing
Data}, 3rd ed. Wiley.\\[0.15em]
Ben-David, S. et al. (2010). A theory of learning from different domains.
\emph{Machine Learning}, 79(1--2), 151--175.\\[0.15em]
Kish, L. (1965). \emph{Survey Sampling}. Wiley.\\[0.15em]
Efron, B., Tibshirani, R. J. (1993). \emph{An Introduction to the Bootstrap}.
Chapman \& Hall.
}

\end{document}
